% ---------------------------------------------------------------------------
% A primer on the natural-sequence side of the comparison: what the influenza A
% hemagglutinin subtype groups are, why influenza B is left out, where the
% sequences come from, and what redundancy clustering, DCA-style reweighting and
% MAFFT alignment do. Companion to softmax/monte_carlo/ha_natural_runbook.md
% (the commands) and softmax/monte_carlo/ha_site_entropy.py (the code).
%
% Every factual statement carries a source. Statements with no source are
% measurements made in this repository, and say so.
%
% Build:  cd softmax/docs && pdflatex ha_natural_methods.tex && pdflatex ha_natural_methods.tex
% ---------------------------------------------------------------------------
\documentclass[11pt,a4paper]{article}

\usepackage[T1]{fontenc}
\usepackage[utf8]{inputenc}
\usepackage[margin=2.5cm]{geometry}
\usepackage{amsmath,amssymb}
\usepackage{booktabs}
\usepackage{array}
\usepackage[hidelinks]{hyperref}
\sloppy

\newcommand{\file}[1]{\texttt{#1}}
\newcommand{\Neff}{N_{\mathrm{eff}}}

\title{The natural side of the comparison\\
       \large Influenza A hemagglutinin subtypes, sequence databases,\\
       redundancy weighting and alignment}
\author{}
\date{3 November 2026}

\begin{document}
\maketitle

\begin{abstract}
\noindent
The sampler produces a per-site entropy over the sequences it accepts. To say
whether that entropy means anything we compare it with the per-site entropy of
real hemagglutinin (HA) sequences. Until now the natural side was 19\,380 human
H3N2 sequences; this note explains what we are replacing that with and why, and
defines every piece of jargon in the plan: subtype groups, influenza B, the
NCBI, BV-BRC and GISAID databases, CD-HIT and MMseqs2 clustering, DCA-style
sequence weights, and MAFFT alignment. Section~\ref{sec:weighting} is the long
version of the one real judgement call in the pipeline.
\end{abstract}

\section{Why look past human H3N2}

Our reference sequence, \file{zero\_polymer}, is the 566-residue HA0 precursor
of A/New York/392/2004 (H3N2), accession \file{YP\_308839.1}
(\file{softmax/code/tests/references.py}; the length and accession are from
\file{FASTA/H3N2/sequences\_H3N2\_566\_homo\_segment4.csv} in this repository).
So human H3N2 was the obvious natural comparison, and
\file{FASTA/H3N2/site\_entropy\_566.csv} holds it.

The trouble is that human H3N2 is one virus lineage in one host, drifting
antigenically since 1968. A site can be variable there because the fold
tolerates variation, or because antibody pressure actively rewards change at
that site; and a site can be conserved there because the fold needs it, or just
because 50 years is not long enough for anything to have happened. The entropy
of one lineage mixes structural constraint with that lineage's history.

Across many HA subtypes, which have been separated for far longer and sit in
different hosts, a site that is still conserved is conserved because the protein
needs it. That is the quantity the sampler should be reproducing. This is the
standard argument for using deep, diverse alignments in sequence-statistics
work: the covariation and conservation signal that corresponds to physical
constraint emerges from alignments spanning large evolutionary distances
\cite{morcos2011,cocco2018}.

\section{The two subtype groups}
\label{sec:groups}

Influenza A viruses are named by their two surface proteins, HA and
neuraminidase (NA): H3N2 means HA subtype 3, NA subtype 2. Eighteen HA subtypes
are known, H1 to H18, the last two from bats \cite{tong2013}. A subtype is an
antigenic and sequence-divergence category: different subtypes do not
cross-react serologically.

The 18 subtypes fall into exactly two phylogenetic clades, called
\emph{group~1} and \emph{group~2}. The split is old and shows up in every
analysis of HA sequences, going back to the first comparison of all the then
known subtypes \cite{air1981,nobusawa1991}:

\begin{center}
\begin{tabular}{@{}ll@{}}
\toprule
group 1 & H1, H2, H5, H6, H8, H9, H11, H12, H13, H16 (and H17, H18, from bats) \\
group 2 & H3, H4, H7, H10, H14, H15 \\
\bottomrule
\end{tabular}
\end{center}

\noindent
Our H3 reference is in group~2, and H4 and H14 are its closest relatives
\cite{air1981,nobusawa1991}.

\paragraph{What actually differentiates them.}
All HAs have the same architecture: a membrane-distal globular head built from
HA1, which carries the receptor-binding site, on a stalk built mostly from HA2,
which carries the fusion machinery \cite{wilson1981,skehel2000}. The groups are
not different folds; they differ in detail, and the differences that matter are:

\begin{itemize}
  \item \textbf{Sequence divergence.} Between groups, HAs are much further
    apart than within a group. In this repository's own two-subtype test, the
    median identity to the H3 \file{zero\_polymer} reference was 0.966 for human
    H3 sequences and 0.435 for H1 sequences (measured by
    \file{ha\_site\_entropy.py}, which prints it per subtype). Group~1 against
    group~2 is the deepest division inside influenza A HA
    \cite{air1981,nobusawa1991}.
  \item \textbf{Stalk conformation and antibody epitopes.} The clearest
    functional evidence that the group boundary is real comes from broadly
    neutralising antibodies that bind the conserved stalk: CR6261 neutralises
    group~1 HAs but not group~2 \cite{ekiert2009}, CR8020 is the mirror image,
    binding group~2 and not group~1 \cite{ekiert2011}, and antibodies that cover
    both groups are rare enough to have been reported as a result in their own
    right \cite{corti2011}. The two groups present measurably different
    surfaces in the region that is otherwise the most conserved part of the
    molecule.
  \item \textbf{Receptor-binding site details.} Group-specific differences in
    the loops around the receptor-binding site change which sialic acid linkages
    are preferred and how the site can adapt, which is central to the question
    of avian-to-human transmission \cite{russell2004,skehel2000}.
\end{itemize}

\paragraph{Why two files and not one.}
Within group~2, the alignment to \file{zero\_polymer}'s 566 positions is
trustworthy at essentially every position, so a group~2 entropy file is the
conservative measurement. Adding group~1 roughly triples the evolutionary depth,
which is what we want, but the HA1 head alignment gets looser and more columns
become unreliable. Building both and comparing them is cheap, so the plan is to
build both: disagreement between them is itself informative, and agreement means
the result does not hinge on the choice.

\section{Influenza B, and why it is excluded}

Influenza A and influenza B are different genera of the same family,
\emph{Orthomyxoviridae}: \emph{Alphainfluenzavirus} and
\emph{Betainfluenzavirus} \cite{ictv2024}. The practical differences:

\begin{itemize}
  \item \textbf{Host range and reservoir.} Influenza A has a huge wild bird
    reservoir, with many subtypes circulating in many hosts, which is what makes
    its HA so diverse and is the reason it can produce pandemics by reassortment.
    Influenza B is essentially a human virus, with seals the only other
    established host \cite{osterhaus2000}, so it has nothing like influenza A's
    diversity.
  \item \textbf{No subtypes, two lineages.} Influenza B is not divided into HA
    subtypes. It has two antigenically distinct lineages, B/Victoria/2/87-like
    and B/Yamagata/16/88-like, which diverged in the 1980s \cite{rota1990}.
    Compared with A's 18 subtypes that is a very shallow set.
  \item \textbf{No reassortment with A.} The two genera do not exchange
    segments, so their HAs have evolved independently throughout.
  \item \textbf{Same fold, far apart in sequence.} Influenza B HA has the same
    overall structure as influenza A HA, including the HA1/HA2 cleavage and the
    fusion peptide \cite{wang2008}. But it is the most distant member of the
    family from our reference, and this is the problem: a position-by-position
    mapping onto \file{zero\_polymer}'s 566 sites stops being meaningful when
    the alignment itself is uncertain, and the per-site entropy we want is
    exactly a per-position quantity. The round figure usually quoted for A
    against B HA identity is about 25\%, which is in the range where alignment
    columns, not biology, would drive the result. \emph{Treat that 25\% as
    approximate: measure it on your own alignment before quoting it in the
    thesis.}
\end{itemize}

\noindent
So influenza B is excluded not because it is irrelevant but because at that
distance the mapping step, not the conservation signal, would decide the answer.

\section{Where the sequences come from}

Three sources are worth knowing. The code only needs FASTA files, so any of them
works.

\paragraph{NCBI (the Influenza Virus Resource, now NCBI Virus).}
The one you already know. It was published as a dedicated influenza resource,
with its own sequence set, annotation pipeline and alignment tools
\cite{bao2008}, and has since been folded into the broader NCBI Virus / Virus
Variation platform \cite{hatcher2017}. It is public, needs no account, and has a
programmatic interface, the E-utilities \cite{eutils}, which is what
\file{ha\_site\_entropy.py fetch} uses: one query per subtype against the
protein database, selecting by organism (the taxonomy has a node per HxNy
subtype), by protein name, and by sequence length.

\paragraph{BV-BRC (formerly FluDB / IRD).}
The Influenza Research Database (IRD), whose website was \file{fludb.org}, was a
US-funded influenza-specific resource: sequences plus curated metadata (host,
location, date, passage history) and analysis tools, built for surveillance work
\cite{squires2012}. In 2019--2022 IRD was merged with PATRIC (bacteria) and ViPR
(other viruses) into a single centre, the Bacterial and Viral Bioinformatics
Resource Center, BV-BRC, at \file{bv-brc.org} \cite{olson2023}. So
``FluDB''/``IRD''/``BV-BRC'' are three names for the same lineage of resource.
For us its advantage over NCBI is the curated metadata: if you want ``avian H3
only'' or ``exclude egg-passaged isolates'', BV-BRC makes that a filter rather
than a text search.

\paragraph{GISAID.}
The Global Initiative on Sharing All Influenza Data, whose influenza database is
EpiFlu \cite{shu2017,elbe2017}. It holds sequences that are not in NCBI,
because many national laboratories deposit there under its terms. Those terms
are the catch: access requires registration and agreement to a data access
agreement that forbids redistribution and requires acknowledging the submitting
laboratories. Practically, that means you may use GISAID data in your analysis
but you cannot commit the FASTA files to this repository, and the acknowledgement
table has to go in the thesis. For our purpose, which is deep subtype coverage
rather than the latest circulating strains, NCBI alone is very likely enough, so
I would only go to GISAID if some subtype turns out to be thinly covered.

\section{Redundancy: clustering and weighting}
\label{sec:redundancy}

\paragraph{The problem.}
Sequence databases are not a sample of the protein's sequence space, they are a
sample of what people sequenced. Human H3N2 and H1N1 have been sequenced tens of
thousands of times because they infect us; H14 has barely been sequenced at all.
Shannon entropy computed on raw counts would therefore mostly describe
surveillance effort. Worse, near-identical sequences make every site look
conserved: if 10\,000 of your 11\,000 sequences are the same recent human
H3N2 clade, the entropy you measure is that clade's internal diversity.

There are two standard cures.

\paragraph{Clustering: CD-HIT and MMseqs2.}
Both are programs that take a set of sequences and group them so that no two
groups are more similar than a chosen identity threshold, keeping one
representative per group. CD-HIT does this with a greedy incremental algorithm:
sort by length, take the longest sequence as the first representative, and
compare each next sequence against existing representatives, making it a new
representative only if it is below the threshold against all of them; short-word
filters make the comparisons fast enough to be practical
\cite{li2006,fu2012}. MMseqs2 is a newer and much faster suite doing the same
job (and sensitive search besides), with a linear-time clustering mode for very
large sets \cite{steinegger2017,steinegger2018}. Clustering at 90--95\% identity
is the usual recipe for ``remove near-duplicates''. The cost is that it throws
information away: a cluster of 10\,000 sequences and a cluster of 2 both end up
as one sequence.

\paragraph{DCA-style weights.}
The alternative keeps every sequence but downweights the crowded ones. It comes
from direct coupling analysis, the Potts-model approach to inferring contacts
and fitness from sequence alignments \cite{morcos2011,cocco2018}. Each sequence
$a$ gets the weight
\begin{equation}
  w_a \;=\; \frac{1}{m_a},
  \qquad
  m_a \;=\; \#\{\,b : \mathrm{identity}(a,b) \ge \theta \,\},
  \label{eq:weights}
\end{equation}
where $b$ runs over all sequences including $a$ itself, so $m_a \ge 1$ and
$w_a \le 1$. A sequence with no close neighbours counts fully; each member of a
cluster of 100 near-identical sequences counts 1/100, so the cluster as a whole
counts about~1. The effective number of sequences is the total weight,
\begin{equation}
  \Neff \;=\; \sum_a w_a ,
  \label{eq:neff}
\end{equation}
and it is the honest sample size: it says roughly how many independent sequences
the alignment is worth. The usual threshold is $\theta = 0.8$, i.e.\ 80\%
identity; Morcos \emph{et al.} use 80\%, define the weight as the inverse count
of sequences above that identity, and define $\Neff$ as the sum of the weights,
reporting only a weak dependence on the precise threshold \cite{morcos2011}.
This is what \file{ha\_site\_entropy.py} implements, with $\theta$ exposed as
\file{--weight-identity}; frequencies and entropies are then computed with $w_a$
in place of a count of~1.

Weighting is the better choice here, and it is what the code does: it is a
smooth correction rather than a hard cut, it keeps every sequence, and it hands
you $\Neff$ as a number you can report, which clustering does not. It is worth
knowing that the value of phylogenetic reweighting has been questioned for
contact prediction specifically \cite{hockenberry2019}; for our purpose, where
the point is precisely not to let surveillance effort set the entropy, the
correction is doing something we want regardless.

\section{Alignment: what MAFFT is for}
\label{sec:mafft}

\paragraph{Why alignment is needed at all.}
The H3N2 file needed none: all 19\,380 sequences were already 566 residues long,
so position 7 meant the same thing in all of them. That stops being true across
subtypes. HA sequences differ in length, with insertions and deletions, so the
7th letter of an H5 sequence is not the 7th site of our reference. A multiple
sequence alignment inserts gaps so that each column holds positions that are
homologous, i.e.\ descended from the same ancestral position. Only then is a
per-column entropy meaningful.

\paragraph{MAFFT.}
MAFFT is one of the standard multiple sequence alignment programs, originally
built around a fast Fourier transform trick for finding homologous segments
quickly \cite{katoh2002}, and in current use as version 7 with a range of
accuracy/speed settings \cite{katoh2013}. It is free, fast enough for thousands
of HA sequences on a laptop (8 seconds for 3\,165 sequences in our test), and
scriptable. Any comparable program (MUSCLE, Clustal Omega, HMMER's
\file{hmmalign} against the Pfam PF00509 model) would do the same job.

\paragraph{What we align, and to what.}
All the kept sequences \emph{together with} the \file{zero\_polymer} reference,
in one run (\file{--align msa}). Then, and this is the important step, we keep
only the alignment columns where the reference has a residue rather than a gap.
That projects the alignment onto the reference's own numbering: insertions
relative to \file{zero\_polymer} are discarded, deletions appear as gap
characters, and the output has exactly 566 rows, one per site of our reference,
which is what the comparison needs. The code checks that the reference comes out
of the alignment unchanged, and fails loudly otherwise.

The alternative mode (\file{--align add}, MAFFT's \file{--add --keeplength}
\cite{katoh2012}) aligns each sequence against the reference alone. It is faster
and fine within one subtype, which is why it reproduces the existing H3N2 file,
but a full multiple alignment is better across subtypes, because the other
sequences help place the indels.

\paragraph{Gappy columns.}
If a column is mostly gaps, the handful of residues in it carry no reliable
entropy, and treating ``gap'' as a 21st amino acid would be worse: it would
report alignment uncertainty as biological variability. So the code computes the
weighted gap fraction per site and writes \file{NaN} for the entropy wherever it
exceeds \file{--max-gap} (30\% by default), and
\file{structure\_zero\_polymer.py} excludes those sites from the rank
correlation, the overlap counts and the box plots.

\section{The one real judgement call: the weighting threshold}
\label{sec:weighting}

This is the long version of the caveat in \file{ha\_natural\_runbook.md}.

Take equation~\eqref{eq:weights} with $\theta = 0.8$ and apply it to a
cross-subtype set. Within a subtype, sequences are typically far above 80\%
identical to each other: human H3 sequences sit at 0.966 median identity to our
reference, so they are all mutual neighbours. Between subtypes they are far
below it: H1 against H3 is about 0.44. The neighbour sets therefore coincide
almost exactly with the subtypes, so $m_a$ is about the size of $a$'s subtype,
each subtype's weights sum to about~1, and
\begin{equation}
  \Neff \;\approx\; \text{(number of subtypes)} .
\end{equation}
Our two-subtype test came out at $\Neff = 2.0$, which looks alarming until you
see that it is exactly what the formula is supposed to do.

The consequence is what matters. With these weights the entropy at a site is
almost entirely \emph{between}-subtype variation: it asks ``do the different
subtypes agree at this position?'' and is nearly blind to variation inside a
subtype, including all of human H3N2's antigenic drift. That is a defensible and
arguably the right definition of ``what the fold requires'', since the
within-subtype signal is the part most contaminated by host-specific immune
pressure and by sequencing effort. But it is a choice, not a fact, so it should
be shown not to matter:

\begin{itemize}
  \item \file{--weight-identity 0.8} (the default): one effective sequence per
    subtype; pure between-subtype conservation.
  \item \file{--weight-identity 0.95}: within-subtype drift is kept, since
    sequences only count as neighbours when they are nearly identical, and each
    subtype contributes in proportion to how much of it was sequenced. A middle
    course, and closer to what the original H3N2-only file measured.
  \item \file{--weight-identity 0} (no weighting): raw counts, so the entropy is
    dominated by whichever subtypes have been sequenced most, in practice H1, H3
    and H5. This is the sampling-biased extreme and is useful mainly as a
    control.
\end{itemize}

\noindent
Build the file at 0.8 and at 0.95, run the K-site comparison on both, and report
that the conclusion is the same. If it is not the same, that is a finding about
which kind of conservation the sampler is reproducing, and more interesting than
either number alone.

\paragraph{Correction, from the first real run.}
The paragraph above says $\Neff$ comes out near the number of subtypes at
$\theta = 0.8$, with each subtype contributing about~1. On the real 16-subtype
data that is only true for the subtypes that have no close sister. Sister
subtypes are themselves more than 80\% identical, so they merge into one
cluster, and a small subtype next to a large one is swallowed: H15 (11 unique
sequences, sister to H7 with 1\,712) came out with $\Neff = 0.0$, and H13 and
H16 with 0.7 and 0.4. At $\theta = 0.8$ the entropy is therefore one vote per
\emph{cluster of sister subtypes}, not per subtype, and the rarest subtypes drop
out of the average altogether. That is a worse definition than intended, so
prefer the $\theta = 0.95$ files and keep 0.8 only as a comparison.

\section{What the answer should look like}

A prediction worth writing down before looking, because it makes the result
falsifiable. If the sampler is finding real structural constraint, then in the
natural alignment:

\begin{itemize}
  \item HA2 should be conserved: the fusion peptide at its N-terminus, the long
    central helix, the stalk. This region does the membrane fusion and is
    sterically constrained, and it is where the broadly neutralising antibodies
    bind precisely because it cannot change much
    \cite{skehel2000,ekiert2009,ekiert2011}.
  \item The HA1 head should be variable, especially the five antigenic sites
    around the receptor-binding site that were mapped from the original
    structure and from drift variants \cite{wiley1981,wilson1981}.
\end{itemize}

\noindent
In the repository's two-subtype test this split is already visible: median
entropy 1.00 bits over HA1 against 0.084 bits over HA2 (measured by
\file{ha\_site\_entropy.py}; the HA1/HA2 boundary is found from the fusion
peptide motif G[IL]FGAIAGF).

So if the T12 K-sites concentrate in HA2 and the stalk, that is evidence the
sampler is doing something physical, before any correlation coefficient is
computed. And if they scatter uniformly, or sit in the antigenic sites, no
Spearman value will rescue the result.

\begin{thebibliography}{99}

\bibitem{air1981} G.~M. Air,
  \emph{Sequence relationships among the hemagglutinin genes of 12 subtypes of
  influenza A virus}, Proc. Natl. Acad. Sci. USA \textbf{78}, 7639--7643 (1981).

\bibitem{bao2008} Y.~Bao, P.~Bolotov, D.~Dernovoy, B.~Kiryutin, L.~Zaslavsky,
  T.~Tatusova, J.~Ostell, D.~Lipman,
  \emph{The influenza virus resource at the National Center for Biotechnology
  Information}, J. Virol. \textbf{82}, 596--601 (2008).

\bibitem{cocco2018} S.~Cocco, C.~Feinauer, M.~Figliuzzi, R.~Monasson,
  M.~Weigt,
  \emph{Inverse statistical physics of protein sequences: a key issues review},
  Rep. Prog. Phys. \textbf{81}, 032601 (2018).

\bibitem{corti2011} D.~Corti \emph{et al.},
  \emph{A neutralizing antibody selected from plasma cells that binds to group 1
  and group 2 influenza A hemagglutinins}, Science \textbf{333}, 850--856
  (2011).

\bibitem{ekiert2009} D.~C. Ekiert, G.~Bhabha, M.-A. Elsliger, R.~H.~E. Friesen,
  M.~Jongeneelen, M.~Throsby, J.~Goudsmit, I.~A. Wilson,
  \emph{Antibody recognition of a highly conserved influenza virus epitope},
  Science \textbf{324}, 246--251 (2009).

\bibitem{ekiert2011} D.~C. Ekiert \emph{et al.},
  \emph{A highly conserved neutralizing epitope on group 2 influenza A
  viruses}, Science \textbf{333}, 843--850 (2011).

\bibitem{elbe2017} S.~Elbe, G.~Buckland-Merrett,
  \emph{Data, disease and diplomacy: GISAID's innovative contribution to global
  health}, Global Challenges \textbf{1}, 33--46 (2017).

\bibitem{eutils} E.~Sayers,
  \emph{A General Introduction to the E-utilities}, Entrez Programming
  Utilities Help, NCBI (updated continuously).
  \url{https://www.ncbi.nlm.nih.gov/books/NBK25497/}

\bibitem{fu2012} L.~Fu, B.~Niu, Z.~Zhu, S.~Wu, W.~Li,
  \emph{CD-HIT: accelerated for clustering the next-generation sequencing
  data}, Bioinformatics \textbf{28}, 3150--3152 (2012).

\bibitem{hatcher2017} E.~L. Hatcher \emph{et al.},
  \emph{Virus Variation Resource -- improved response to emergent viral
  outbreaks}, Nucleic Acids Res. \textbf{45}, D482--D490 (2017).

\bibitem{hockenberry2019} A.~J. Hockenberry, C.~O. Wilke,
  \emph{Phylogenetic weighting does little to improve the accuracy of
  evolutionary coupling analyses}, Entropy \textbf{21}, 1000 (2019).

\bibitem{ictv2024} International Committee on Taxonomy of Viruses,
  \emph{Orthomyxoviridae}, ICTV Report on Virus Taxonomy.
  \url{https://ictv.global/report/chapter/orthomyxoviridae}

\bibitem{katoh2002} K.~Katoh, K.~Misawa, K.~Kuma, T.~Miyata,
  \emph{MAFFT: a novel method for rapid multiple sequence alignment based on
  fast Fourier transform}, Nucleic Acids Res. \textbf{30}, 3059--3066 (2002).

\bibitem{katoh2012} K.~Katoh, M.~C. Frith,
  \emph{Adding unaligned sequences into an existing alignment using MAFFT and
  LAST}, Bioinformatics \textbf{28}, 3144--3146 (2012).

\bibitem{katoh2013} K.~Katoh, D.~M. Standley,
  \emph{MAFFT multiple sequence alignment software version 7: improvements in
  performance and usability}, Mol. Biol. Evol. \textbf{30}, 772--780 (2013).

\bibitem{li2006} W.~Li, A.~Godzik,
  \emph{Cd-hit: a fast program for clustering and comparing large sets of
  protein or nucleotide sequences}, Bioinformatics \textbf{22}, 1658--1659
  (2006).

\bibitem{morcos2011} F.~Morcos, A.~Pagnani, B.~Lunt, A.~Bertolino,
  D.~S. Marks, C.~Sander, R.~Zecchina, J.~N. Onuchic, T.~Hwa, M.~Weigt,
  \emph{Direct-coupling analysis of residue coevolution captures native contacts
  across many protein families}, Proc. Natl. Acad. Sci. USA \textbf{108},
  E1293--E1301 (2011).
  [Reweighting: 80\% identity threshold, weight $1/m_a$, $\Neff$ as the sum of
  the weights, with only a weak dependence on the threshold.]

\bibitem{nobusawa1991} E.~Nobusawa, T.~Aoyama, H.~Kato, Y.~Suzuki,
  Y.~Tateno, K.~Nakajima,
  \emph{Comparison of complete amino acid sequences and receptor-binding
  properties among 13 serotypes of hemagglutinins of influenza A viruses},
  Virology \textbf{182}, 475--485 (1991).

\bibitem{olson2023} R.~D. Olson \emph{et al.},
  \emph{Introducing the Bacterial and Viral Bioinformatics Resource Center
  (BV-BRC): a resource combining PATRIC, IRD and ViPR}, Nucleic Acids Res.
  \textbf{51}, D678--D689 (2023).

\bibitem{osterhaus2000} A.~D.~M.~E. Osterhaus, G.~F. Rimmelzwaan,
  B.~E.~E. Martina, T.~M. Bestebroer, R.~A.~M. Fouchier,
  \emph{Influenza B virus in seals}, Science \textbf{288}, 1051--1053 (2000).

\bibitem{rota1990} P.~A. Rota, T.~R. Wallis, M.~W. Harmon, J.~S. Rota,
  A.~P. Kendal, K.~Nerome,
  \emph{Cocirculation of two distinct evolutionary lineages of influenza type B
  virus since 1983}, Virology \textbf{175}, 59--68 (1990).

\bibitem{russell2004} R.~J. Russell \emph{et al.},
  \emph{H1 and H7 influenza haemagglutinin structures extend a structural
  classification of haemagglutinin subtypes}, Virology \textbf{325}, 287--296
  (2004).

\bibitem{shu2017} Y.~Shu, J.~McCauley,
  \emph{GISAID: Global initiative on sharing all influenza data -- from vision
  to reality}, Eurosurveillance \textbf{22}, 30494 (2017).

\bibitem{skehel2000} J.~J. Skehel, D.~C. Wiley,
  \emph{Receptor binding and membrane fusion in virus entry: the influenza
  hemagglutinin}, Annu. Rev. Biochem. \textbf{69}, 531--569 (2000).

\bibitem{squires2012} R.~B. Squires \emph{et al.},
  \emph{Influenza Research Database: an integrated bioinformatics resource for
  influenza research and surveillance}, Influenza Other Respir. Viruses
  \textbf{6}, 404--416 (2012).

\bibitem{steinegger2017} M.~Steinegger, J.~S\"oding,
  \emph{MMseqs2 enables sensitive protein sequence searching for the analysis of
  massive data sets}, Nat. Biotechnol. \textbf{35}, 1026--1028 (2017).

\bibitem{steinegger2018} M.~Steinegger, J.~S\"oding,
  \emph{Clustering huge protein sequence sets in linear time}, Nat. Commun.
  \textbf{9}, 2542 (2018).

\bibitem{tong2013} S.~Tong \emph{et al.},
  \emph{New world bats harbor diverse influenza A viruses}, PLoS Pathog.
  \textbf{9}, e1003657 (2013).

\bibitem{wang2008} Q.~Wang, F.~Tian, X.~Feng, H.~Cheng, J.~Yan, G.~F. Gao,
  \emph{Crystal structure of unliganded influenza B virus hemagglutinin},
  J. Virol. \textbf{82}, 3011--3020 (2008).

\bibitem{wiley1981} D.~C. Wiley, I.~A. Wilson, J.~J. Skehel,
  \emph{Structural identification of the antibody-binding sites of Hong Kong
  influenza haemagglutinin and their involvement in antigenic variation},
  Nature \textbf{289}, 373--378 (1981).

\bibitem{wilson1981} I.~A. Wilson, J.~J. Skehel, D.~C. Wiley,
  \emph{Structure of the haemagglutinin membrane glycoprotein of influenza virus
  at 3\,\AA{} resolution}, Nature \textbf{289}, 366--373 (1981).

\end{thebibliography}

\end{document}
